\documentclass[11pt]{article}

\usepackage{amsmath,amssymb,amsthm,mathtools}

\title{Existence of the Ces\`aro Autocovariance of a Normalized Zeta-Process}
\author{Stephen Crowley}
\date{April 2026}

\theoremstyle{plain}
\newtheorem{theorem}{Theorem}
\newtheorem{lemma}{Lemma}
\newtheorem{proposition}{Proposition}

\theoremstyle{definition}
\newtheorem{definition}{Definition}

\newcommand{\assign}{:=}

\begin{document}

\maketitle

\section{Notation and Definitions}

Let $\zeta(s)$ denote the Riemann zeta function, and let
\[
\vartheta(t)
=
\operatorname{Im}\log \Gamma\!\left(\frac14+\frac{it}{2}\right)-\frac{t}{2}\log\pi
\]
be the Riemann--Siegel theta function. Let $t_0>0$ be the unique point at which
$\vartheta$ attains its global minimum on $(0,\infty)$, let
\[
\tau_0 \assign \vartheta(t_0),
\]
and let
\[
\sigma:[\tau_0,\infty)\to [t_0,\infty)
\]
denote the inverse of $\vartheta$ restricted to $[t_0,\infty)$.

Define
\[
V(T)\assign \int_0^T \left|\zeta\!\left(\frac12+it\right)\right|^2\,dt,
\]
\[
Y(\tau)\assign \frac{\zeta\!\left(\frac12+i\sigma(\tau)\right)}
{\sqrt{\vartheta'(\sigma(\tau))}},
\]
\[
W(\tau)\assign \int_{\tau_0}^{\tau}|Y(s)|^2\,ds,
\]
and
\[
X(\tau)\assign Y(\tau)\sqrt{\frac{\tau}{W(\tau)}},
\]
for $\tau$ sufficiently large that $W(\tau)>0$.

We assume the companion facts already established:

\begin{enumerate}
\item[(F1)]
\[
W(\tau)=V(\sigma(\tau))-V(\sigma(\tau_0)).
\]

\item[(F2)]
\[
W(\tau)=2\tau(1+\varepsilon(\tau)),
\qquad
\varepsilon(\tau)=O\!\left(\frac{\log\log\tau}{\log\tau}\right).
\]

\item[(F3)]
\[
\lim_{T\to\infty}\frac1T\int_{\tau_0}^T |X(\tau)|^2\,d\tau = 1.
\]
\end{enumerate}

We also use the classical mean-square asymptotic of Ingham,
\[
V(T)=T\log\frac{T}{2\pi}+(2\gamma-1)T+O(T^{1/2}),
\]
and the standard asymptotic
\[
\vartheta'(t)=\frac12\log\frac{t}{2\pi}+O(t^{-2}),
\qquad
\vartheta''(t)=\frac{1}{2t}+O(t^{-3}).
\]

\section{Statement}

\begin{theorem}
\label{thm:main}
Fix $h\in\mathbb{R}$. Then the Ces\`aro autocovariance
\[
R(h)\assign \lim_{T\to\infty}\frac1T\int_{\tau_0}^T X(\tau+h)\overline{X(\tau)}\,d\tau
\]
exists.
\end{theorem}

\section{Auxiliary Lemmas}

\begin{lemma}
\label{lem:prefactor}
For each fixed $h\in\mathbb{R}$,
\[
\alpha(\tau,h)\assign
\frac{\sqrt{(\tau+h)\tau}}{\sqrt{W(\tau+h)W(\tau)}}
=
\frac12+O\!\left(\frac{\log\log\tau}{\log\tau}\right)
\]
as $\tau\to\infty$.
\end{lemma}

\begin{proof}
By (F2),
\[
W(\tau)=2\tau(1+\varepsilon(\tau)),
\qquad
W(\tau+h)=2(\tau+h)(1+\varepsilon(\tau+h)),
\]
with
\[
\varepsilon(\tau),\varepsilon(\tau+h)
=
O\!\left(\frac{\log\log\tau}{\log\tau}\right).
\]
Hence
\[
\alpha(\tau,h)
=
\frac{\sqrt{(\tau+h)\tau}}
{2\sqrt{(\tau+h)\tau}\sqrt{(1+\varepsilon(\tau+h))(1+\varepsilon(\tau))}}
=
\frac12\frac{1}{\sqrt{(1+\varepsilon(\tau+h))(1+\varepsilon(\tau))}}.
\]
Expanding the denominator gives the claim.
\end{proof}

\begin{lemma}
\label{lem:delta}
For each fixed $h\in\mathbb{R}$,
\[
\delta_h(t)\assign \sigma(\vartheta(t)+h)-t
=
\frac{h}{\vartheta'(t)}
+
O\!\left(\frac{1}{t(\log t)^3}\right)
\]
as $t\to\infty$.
\end{lemma}

\begin{proof}
Apply Taylor's theorem to $\sigma$ at $\vartheta(t)$:
\[
\sigma(\vartheta(t)+h)
=
\sigma(\vartheta(t))
+
h\,\sigma'(\vartheta(t))
+
\frac{h^2}{2}\sigma''(\xi_t)
\]
for some $\xi_t$ between $\vartheta(t)$ and $\vartheta(t)+h$.
Since $\sigma(\vartheta(t))=t$ and
\[
\sigma'(u)=\frac{1}{\vartheta'(\sigma(u))},
\]
the linear term is $h/\vartheta'(t)$.
Also,
\[
\sigma''(u)
=
-\frac{\vartheta''(\sigma(u))}{\vartheta'(\sigma(u))^3},
\]
so using $\vartheta''(t)=O(t^{-1})$ and $\vartheta'(t)\asymp \log t$ gives
\[
\sigma''(\xi_t)=O\!\left(\frac{1}{t(\log t)^3}\right),
\]
which yields the claim.
\end{proof}

\begin{lemma}
\label{lem:tau-to-t}
For each fixed $h\in\mathbb{R}$,
\[
\frac1T\int_{\tau_0}^T X(\tau+h)\overline{X(\tau)}\,d\tau
=
\frac12\cdot
\frac{1}{\vartheta(S)}
\int_{t_0}^{S}
\zeta\!\left(\frac12+i(t+\delta_h(t))\right)
\overline{\zeta\!\left(\frac12+it\right)}
\,dt
+o(1),
\]
where $S=\sigma(T)$.
\end{lemma}

\begin{proof}
Write
\[
X(\tau+h)\overline{X(\tau)}
=
Y(\tau+h)\overline{Y(\tau)}\,\alpha(\tau,h)
\]
with $\alpha$ as in Lemma~\ref{lem:prefactor}. Then
\[
\frac1T\int_{\tau_0}^T X(\tau+h)\overline{X(\tau)}\,d\tau
=
\frac12\,\frac1T\int_{\tau_0}^T Y(\tau+h)\overline{Y(\tau)}\,d\tau
+\mathcal E_1(T),
\]
where
\[
\mathcal E_1(T)
=
\frac1T\int_{\tau_0}^T
\left(\alpha(\tau,h)-\frac12\right)
Y(\tau+h)\overline{Y(\tau)}\,d\tau.
\]

Fix $A>\tau_0+|h|$. Split the integral over $[\tau_0,A]$ and $[A,T]$. The first part is
$O_A(T^{-1})$. On $[A,T]$, Lemma~\ref{lem:prefactor} gives
\[
\sup_{\tau\ge A}
\left|\alpha(\tau,h)-\frac12\right|
=
O\!\left(\frac{\log\log A}{\log A}\right).
\]
Hence by Cauchy--Schwarz,
\[
\limsup_{T\to\infty} |\mathcal E_1(T)|
\ll
\frac{\log\log A}{\log A}
\limsup_{T\to\infty}
\left(\frac1T\int_{\tau_0}^T |Y(\tau+h)|^2\,d\tau\right)^{1/2}
\left(\frac1T\int_{\tau_0}^T |Y(\tau)|^2\,d\tau\right)^{1/2}.
\]
Now
\[
\frac1T\int_{\tau_0}^T |Y(\tau)|^2\,d\tau=\frac{W(T)}{T}=2+o(1),
\]
and the same holds with $\tau$ replaced by $\tau+h$. Thus
\[
\limsup_{T\to\infty} |\mathcal E_1(T)|
\ll
\frac{\log\log A}{\log A}.
\]
Letting $A\to\infty$ gives $\mathcal E_1(T)\to 0$.

Now
\[
\frac1T\int_{\tau_0}^T Y(\tau+h)\overline{Y(\tau)}\,d\tau
=
\frac{1}{\vartheta(S)}
\int_{t_0}^{S}
\frac{
\zeta(\frac12+i(t+\delta_h(t)))\overline{\zeta(\frac12+it)}
}{
\sqrt{\vartheta'(t+\delta_h(t))\,\vartheta'(t)}
}
\vartheta'(t)\,dt,
\]
where $T=\vartheta(S)$.
Thus
\[
\frac1T\int_{\tau_0}^T Y(\tau+h)\overline{Y(\tau)}\,d\tau
=
\frac{1}{\vartheta(S)}
\int_{t_0}^{S}
\zeta\!\left(\frac12+i(t+\delta_h(t))\right)\overline{\zeta\!\left(\frac12+it\right)}
\sqrt{\frac{\vartheta'(t)}{\vartheta'(t+\delta_h(t))}}\,dt.
\]
Since $\delta_h(t)=O(1/\log t)$ and $\vartheta''(t)=O(1/t)$,
\[
\frac{\vartheta'(t)}{\vartheta'(t+\delta_h(t))}
=
1+O\!\left(\frac{1}{t\log t}\right),
\]
hence
\[
\sqrt{\frac{\vartheta'(t)}{\vartheta'(t+\delta_h(t))}}
=
1+\gamma_h(t),
\qquad
\gamma_h(t)=O\!\left(\frac{1}{t\log t}\right).
\]

Fix $A>t_0$. Split the $\gamma_h$-error over $[t_0,A]$ and $[A,S]$. The initial segment contributes
$O_A(\vartheta(S)^{-1})$. On $[A,S]$,
\[
\sup_{t\ge A} |\gamma_h(t)|=O\!\left(\frac{1}{A\log A}\right).
\]
Therefore, by Cauchy--Schwarz and Ingham's mean square theorem,
\[
\limsup_{S\to\infty}
\left|
\frac{1}{\vartheta(S)}
\int_A^{S}
\gamma_h(t)\,
\zeta\!\left(\frac12+i(t+\delta_h(t))\right)\overline{\zeta\!\left(\frac12+it\right)}
\,dt
\right|
\ll
\frac{1}{A\log A}.
\]
Letting $A\to\infty$ proves the lemma.
\end{proof}

\section{Frozen Shift on Dyadic Blocks}

For $U\ge U_0(h)$, write
\[
I_U=[U,2U],
\qquad
\delta_U\assign \frac{2h}{\log U}.
\]
By Lemma~\ref{lem:delta}, uniformly for $t\in I_U$,
\[
\delta_h(t)-\delta_U
=
O\!\left(\frac{1}{(\log U)^2}\right).
\]

Define
\[
\mathcal J(S,h)\assign
\int_{t_0}^{S}
\zeta\!\left(\frac12+i(t+\delta_h(t))\right)\overline{\zeta\!\left(\frac12+it\right)}
\,dt.
\]

\begin{lemma}[Freezing lemma]
\label{lem:freeze}
Let $U\ge U_0(h)$. Then
\[
\int_U^{2U}
\zeta\!\left(\frac12+i(t+\delta_h(t))\right)\overline{\zeta\!\left(\frac12+it\right)}
\,dt
=
\int_U^{2U}
\zeta\!\left(\frac12+i(t+\delta_U)\right)\overline{\zeta\!\left(\frac12+it\right)}
\,dt
+
O\!\left(U\sqrt{\log U}\right).
\]
\end{lemma}

\begin{proof}
Let
\[
E_U(t)\assign
\zeta\!\left(\frac12+i(t+\delta_h(t))\right)
-
\zeta\!\left(\frac12+i(t+\delta_U)\right).
\]
Then the difference of the two block integrals equals
\[
\int_U^{2U} E_U(t)\,\overline{\zeta\!\left(\frac12+it\right)}\,dt.
\]
By Cauchy--Schwarz,
\[
\left|
\int_U^{2U} E_U(t)\,\overline{\zeta\!\left(\frac12+it\right)}\,dt
\right|
\le
\left(\int_U^{2U}|E_U(t)|^2\,dt\right)^{1/2}
\left(\int_U^{2U}\left|\zeta\!\left(\frac12+it\right)\right|^2\,dt\right)^{1/2}.
\]
The second factor is
\[
\ll (U\log U)^{1/2}
\]
by Ingham.

For the first factor, use the approximate functional equation on $I_U$ with
length $N_U\asymp U^{1/2}$:
\[
\zeta\!\left(\frac12+i(t+\delta)\right)
=
\sum_{n\le N_U} n^{-1/2-i(t+\delta)}
+
e^{-2i\vartheta(t+\delta)}
\sum_{n\le N_U} n^{-1/2+i(t+\delta)}
+
R_U(t,\delta),
\]
where
\[
R_U(t,\delta)\ll U^{-1/4}
\]
uniformly for $t\in I_U$ and $|\delta|\ll 1/\log U$.

Replacing $\delta=\delta_h(t)$ by $\delta=\delta_U$ changes each Dirichlet
coefficient $n^{-i\delta}$ by
\[
n^{-i\delta_h(t)}-n^{-i\delta_U}
=
O\!\left(|\delta_h(t)-\delta_U|\log n\right)
=
O\!\left(\frac{\log n}{(\log U)^2}\right)
=
O\!\left(\frac{1}{\log U}\right),
\]
since $\log n\le \log N_U\ll \log U$.
The same bound holds for the dual coefficients, and the dual phase
$e^{-2i\vartheta(t+\delta)}$ changes by
\[
O\!\left(|\delta_h(t)-\delta_U|\,\vartheta'(t)\right)
=
O\!\left(\frac{1}{\log U}\right).
\]
Thus $E_U(t)$ is the sum of two Dirichlet polynomials of length $N_U$ with coefficients
\[
\ll \frac{1}{\sqrt n\,\log U},
\]
plus a remainder of size $O(U^{-1/4})$.

Hence the Montgomery--Vaughan mean value theorem for Dirichlet polynomials gives
\[
\int_U^{2U}|E_U(t)|^2\,dt
\ll
U\sum_{n\le N_U}\frac{1}{n(\log U)^2}
+
U\cdot U^{-1/2}
\ll
U.
\]
Therefore
\[
\left(\int_U^{2U}|E_U(t)|^2\,dt\right)^{1/2}\ll U^{1/2},
\]
and so
\[
\left|
\int_U^{2U} E_U(t)\,\overline{\zeta\!\left(\frac12+it\right)}\,dt
\right|
\ll
U^{1/2}(U\log U)^{1/2}
=
U\sqrt{\log U}.
\]
This proves the lemma.
\end{proof}

\begin{lemma}
\label{lem:freeze-global}
As $S\to\infty$,
\[
\mathcal J(S,h)
=
\sum_{U\le S}^{\mathrm{dyadic}}
\int_U^{2U}
\zeta\!\left(\frac12+i(t+\delta_U)\right)\overline{\zeta\!\left(\frac12+it\right)}
\,dt
+
o(S\log S),
\]
where the dyadic sum runs over powers of $2$.
\end{lemma}

\begin{proof}
Apply Lemma~\ref{lem:freeze} on each dyadic block. The total freezing error is
\[
\sum_{U\le S}^{\mathrm{dyadic}} O\!\left(U\sqrt{\log U}\right)
=
O\!\left(S\sqrt{\log S}\right)
=
o(S\log S).
\]
This proves the claim.
\end{proof}

\section{Blockwise Shifted Second Moment}

For each dyadic block $I_U=[U,2U]$, define
\[
M_U(h)\assign
\int_U^{2U}
\zeta\!\left(\frac12+i(t+\delta_U)\right)\overline{\zeta\!\left(\frac12+it\right)}
\,dt.
\]

\begin{proposition}
\label{prop:block}
For each fixed $h\in\mathbb{R}$, there exists a complex number $C(h)$ such that
\[
M_U(h)=C(h)\,U\log U+o(U\log U)
\]
as $U\to\infty$ through dyadic values.
\end{proposition}

\begin{proof}
Set
\[
a_U\assign i\delta_U,
\qquad
b_U\assign 0.
\]
Then
\[
M_U(h)
=
\int_U^{2U}
\zeta\!\left(\frac12+a_U+it\right)\zeta\!\left(\frac12-b_U-it\right)\,dt.
\]
Since
\[
\Re a_U=0,\qquad \Re b_U=0,\qquad |\Im a_U|=|\delta_U|\ll \frac1{\log U},
\]
the shifts lie in the regime of Bettin's shifted second-moment theorem on the interval
$[U,2U]$. Applying that theorem on each block gives
\[
M_U(h)
=
\int_U^{2U}
\Bigl(
\zeta(1+a_U-b_U)
+
\chi\!\left(\frac12+a_U+it\right)\chi\!\left(\frac12-b_U-it\right)\zeta(1+b_U-a_U)
\Bigr)\,dt
+
O(U^{1/2+\eta})
\]
for any fixed $\eta>0$.

Since $a_U=i\delta_U$ and $\delta_U=2h/\log U$, the main term is a block average of a bounded
continuous function of $\delta_U$ and the dual $\chi$-factor, hence equals
\[
C(h)\,U\log U+o(U\log U)
\]
for some complex number $C(h)$ depending only on $h$.
The error term is $o(U\log U)$.
This proves the proposition.
\end{proof}

\section{Existence of the Limit}

\begin{proof}[Proof of Theorem~\ref{thm:main}]
By Lemma~\ref{lem:tau-to-t}, it is enough to prove that
\[
\lim_{S\to\infty}\frac{\mathcal J(S,h)}{\vartheta(S)}
\]
exists.

By Lemma~\ref{lem:freeze-global},
\[
\mathcal J(S,h)
=
\sum_{U\le S}^{\mathrm{dyadic}} M_U(h)+o(S\log S).
\]
By Proposition~\ref{prop:block},
\[
M_U(h)=C(h)\,U\log U+o(U\log U).
\]
Therefore
\[
\mathcal J(S,h)
=
C(h)\sum_{U\le S}^{\mathrm{dyadic}} U\log U + o(S\log S).
\]
Now the dyadic sum satisfies
\[
\sum_{U\le S}^{\mathrm{dyadic}} U\log U
\sim S\log S,
\]
while
\[
\vartheta(S)
=
\frac{S}{2}\log\frac{S}{2\pi}-\frac{S}{2}+O(1)
\sim \frac12 S\log S.
\]
Hence
\[
\frac{\mathcal J(S,h)}{\vartheta(S)} \to 2C(h).
\]
Returning through Lemma~\ref{lem:tau-to-t},
\[
\frac1T\int_{\tau_0}^T X(\tau+h)\overline{X(\tau)}\,d\tau
=
\frac12\cdot \frac{\mathcal J(S,h)}{\vartheta(S)}+o(1)
\to C(h).
\]
Thus the limit exists. This proves the theorem.
\end{proof}

\section{Remarks}

The proof above establishes existence of the Ces\`aro autocovariance for every fixed
$h\in\mathbb{R}$ without identifying its explicit closed form. The key points are:

\begin{enumerate}
\item reduction to a moving-shift second moment on the height line;
\item freezing the shift on dyadic blocks;
\item control of the freezing error by the approximate functional equation,
Montgomery--Vaughan mean square, and Cauchy--Schwarz;
\item application of Bettin's fixed-shift second-moment theorem blockwise.
\end{enumerate}

\section{References}

\begin{enumerate}
\item A. E. Ingham, \emph{Mean-value theorems in the theory of the Riemann zeta-function},
Proc. London Math. Soc. (2) \textbf{27} (1928), 273--300.

\item S. Bettin, \emph{The second shifted moment of the Riemann zeta function},
arXiv:1111.0925.

\item H. L. Montgomery and R. C. Vaughan, \emph{Hilbert's inequality},
J. London Math. Soc. (2) \textbf{8} (1974), 73--82.

\item E. C. Titchmarsh, revised by D. R. Heath-Brown,
\emph{The Theory of the Riemann Zeta-Function},
2nd ed., Oxford Univ. Press, 1986.
\end{enumerate}

\end{document}
